% Supplementary reduction for C.6.  This does not settle the general iff.
\begin{proposition}[A purity reduction for homogeneous Chow motives]
\label{c:chow-unramified-purity-reduction}
Let $K$ be a number field, $v$ a finite place with residue field of
cardinality $q$, and $p$ its residue characteristic.  Let
$M=(X,e,m)$ be a rational Chow motive over $K$, where $X$ is smooth
projective of dimension $d$ and the rational algebraic projector
$e\in\CH^d(X\times_KX)_{\Q}$ acts only in cohomological degree $i$.
Suppose that, for one prime $\ell_0\ne p$, the representation
$M_{\ell_0}$ is unramified at $v$ and the eigenvalues of geometric
Frobenius on it are algebraic numbers all of whose complex conjugates
have one common absolute value.  Then $M_\ell$ is unramified at $v$
for every $\ell\ne p$.

In particular, the asserted independence follows for such a motive if
the weight--monodromy conjecture holds for
$H^i(X_{\overline{K_v}},\Q_{\ell_0})$: its unramified direct summand
$M_{\ell_0}$ is then pure of weight $i-2m$.
\end{proposition}
\begin{proof}
Write $V_\ell=eH^i(X_{\overline{K_v}},\Q_\ell)(m)$.
Saito's Theorem~0.1, in its more general form Theorem~3.1,
gives a rational alternating trace independent of $\ell$ for the
composition of an algebraic correspondence and any Weil element
\citep[Theorems~0.1 and 3.1]{c_saito2003weights}.
The theorem extends from integral to rational correspondences by
linearity.  Taking the correspondence $e$, its action on all other
degrees is zero; since it is idempotent and commutes with the Galois
action, one obtains the exact equality
\[
 \operatorname{Tr}(\sigma\mid V_\ell)
 =(-1)^i q^{-m n(\sigma)}
   \sum_a(-1)^a\operatorname{Tr}
   \bigl(e\circ\sigma\mid H^a(X_{\overline{K_v}},\Q_\ell)\bigr).
\]
Thus $\operatorname{Tr}(\sigma\mid V_\ell)$ is a common rational
number for each $\sigma\in W_{K_v}$.  Apply this also to all positive
powers of $\sigma$.  Newton identities show that the characteristic
polynomial of every such $\sigma$ is in $\Q[T]$ and independent of
$\ell$.  In particular the dimension $r$ is common; if $r=0$ there
is nothing to prove.

Fix a lift $\phi\in W_{K_v}$ of geometric Frobenius.  The local
monodromy theorem supplies, for every $\ell\ne p$, a nilpotent operator
$N_\ell$ and an open inertia subgroup on which the action is
$\exp(t_\ell N_\ell)$; it also gives
\[
 F_\ell N_\ell F_\ell^{-1}=q^{-1}N_\ell,
 \qquad F_\ell=\rho_\ell(\phi).
\]
The common characteristic polynomial of $F_\ell$ has roots of the
single absolute value in the hypothesis.  Over an algebraic closure
of the coefficient field, $N_\ell$ maps a generalized
$\lambda$-eigenspace of $F_\ell$ into its generalized
$q^{-1}\lambda$-eigenspace.  The latter eigenvalue cannot occur:
under any embedding of the common algebraic splitting field into
$\C$, multiplication by $q^{-1}$ changes its absolute value.
Consequently $N_\ell=0$ for every $\ell$.

The image of inertia in each $\rho_\ell$ is therefore finite.  For
each $\tau\in I_{K_v}$ the common trace is
$\operatorname{Tr}(\rho_{\ell_0}(\tau))=r$.  A finite-order matrix
in characteristic zero is semisimple with roots of unity as
eigenvalues.  Their sum can equal the positive integer $r$ only if
all of them are $1$, as follows by embedding their cyclotomic field
into $\C$ and using equality in the triangle inequality.  Hence
$\rho_\ell(\tau)=1$ for every $\tau$ and every $\ell\ne p$.

For the last assertion, the weight--monodromy filtration is functorial
for the algebraic projector.  A pure Weil--Deligne direct summand
whose monodromy is zero is concentrated in monodromy degree zero,
so its Frobenius roots all have weight $i-2m$.
\end{proof}

\begin{remark}[Exact scope of this reduction]
\label{c:chow-purity-scope}
For the full group $H^i(X)$, the purity assumption in the proposition
is Serre's unramified-purity conjecture C7
\citep[\S2.3, printed p.~19-07]{c_serre1970local}.
For a prescribed proper summand it is the corresponding
\emph{motivic extension} of that condition; the literal original C7
must not be silently enlarged.  The named weight--monodromy
conjecture for the ambient cohomology is a sufficient alternative.
The use of Saito's theorem requires an algebraic correspondence.
It cannot be applied to an arbitrary absolute-Hodge or motivated
projector without an additional comparison result.  No converse to
any of these named conjectures is proved here, and this proposition
does not close the general-motive requirement in C.6.
\end{remark}
